% from ega3-7-fr.tex
\oldpage[III]{223}
$X=\operatorname{Spec}(A)$, $\mathscr F=\widetilde M$, où $M$ est un $A$-module
de type fini, $\mathscr J=\widetilde{\mathfrak J}$, où $\mathfrak J$ est un
idéal de type fini de $A$ ; le reste du raisonnement est alors inchangé.

\noindent\textbf{$(\mathbf{Err}_{\mathrm{III}},\,30)$} Dans (I, 9.3.5),
remplacer les lignes 1 à 7 du bas de la p. 173 par le texte suivant :

\emph{Proposition (9.3.5). --- Soient $X$ un préschéma, $\mathscr F$ un
$\mathscr O_X$-Module quasi-cohérent de type fini. Alors il existe un
sous-préschéma fermé $Y$ de $X$, dont l'espace sous-jacent est égal à
$\operatorname{Supp}(\mathscr F)$, et un $\mathscr O_Y$-Module quasi-cohérent de
type fini $\mathscr G$ tels que, si $j:Y\to X$ est l'injection canonique,
$\mathscr F$ soit isomorphe à $j_*(\mathscr G)$.}

Il suffira de montrer que l'Idéal $\mathscr J$ de $\mathscr O_X$,
\emph{annulateur de $\mathscr F$}, est \emph{quasi-cohérent} ; on prendra alors
pour $Y$ le sous-préschéma fermé de $X$ défini par $\mathscr J$ (I, 4.1.2), et
comme $\mathscr J\mathscr F=0$, $\mathscr F$ est un
$(\mathscr O_X/\mathscr J)$-Module et on répondra à la question en prenant
$\mathscr G=j^*(\mathscr F)$. Pour voir que $\mathscr J$ est quasi-cohérent, on
peut (la question étant locale) se borner au cas où
$X=\operatorname{Spec}(A)$, $\mathscr F=\widetilde M$, où $M$ est un $A$-module
engendré par un nombre fini d'éléments $x_i\ (1\leq i\leq r)$ ; l'Idéal
$\mathscr J$ est alors l'intersection des annulateurs des $x_i$. Mais
l'annulateur de $x_i$ est le noyau de l'homomorphisme
$\mathscr O_X\to\mathscr F$ correspondant à l'homomorphisme $s\mapsto sx_i$ de
$A$ dans $M$ ; c'est donc bien un Idéal quasi-cohérent (I, 4.1.1) et toute
intersection finie de tels Idéaux est aussi un Idéal quasi-cohérent (I,
1.3.10).
